\section{Experimental Evaluation of Different Output Models}
\label{sec:experiments}

This chapter presents the experimental evaluation of the proposed models. It discusses the specific characteristics of each model, followed by the evaluation metrics and visual comparisons. The hyperparameter-tuning process used to identify a well-performing configuration for a selected combination of output formulation and loss function is then described. This is followed by a broader comparison of models using different output formulations and loss functions. Finally, the best-performing model is evaluated on a previously unseen test dataset.

All experiments were conducted on a computing cluster using one NVIDIA H100 GPU per training run. The initial training runs completed the full 200 epochs and took approximately 30 hours with a batch size of 10. Early stopping was introduced for subsequent experiments; all of these runs terminated before reaching 100 epochs, substantially reducing their training time.

The models were trained using the AdamW optimizer with a weight decay of \(10^{-8}\). The training and validation sets were generated using a fixed random split with seed 0, ensuring identical data partitions across experiments. The split contains $72{,}000$ training samples and $8{,}000$ validation samples. The constant values are listed in table~\ref{tab:training_configuration}. Some values, especially the last two were changed during hyperparameter training. This is made clear when describing the models. 

\begin{table}[htbp]
\centering
\caption{Training configuration.}
\label{tab:training_configuration}
\begin{tabular}{ll}
\toprule
Parameter & Value \\
\midrule
Optimizer & AdamW \\
Maximum epochs & 200 \\
Weight decay & \(10^{-8}\) \\
Early-stopping patience & 10 or 20 epochs \\
Data-split seed & 0 \\
Batch size & 10 \\
Learning rate & \(10^{-4}\) \\
\bottomrule
\end{tabular}
\end{table}

\subsection{Two-output Model}

\begin{figure}[htbp]
    \centering
    \includegraphics[width=0.95\textwidth]{diagrams/twooutputs/group_01_two_output_objective.pdf}
    \caption[Two-output objective comparison]{Training comparison of the two-output objectives. The L1 and reconstruction objectives optimize intensity agreement and therefore reach low reconstruction errors, whereas the Tversky objective optimizes foreground overlap and produces substantially stronger mask metrics.}
    \label{fig:twooutput-objective-comparison}
\end{figure}

\begin{figure}[htbp]
    \centering
    \includegraphics[width=0.99\textwidth]{diagrams/twooutputs/figure_04_dataset_model_heatmap_two_output.pdf}
    \caption[Two-output dataset performance]{Per-dataset performance for the two-output models. The heatmap shows that the aggregate metrics in Table~\ref{tab:two_output_models} hide a strong dataset dependence: compact, well-separated examples are handled more reliably than datasets with larger, more deformed, or more strongly overlapping spots.}
    \label{fig:twooutput-per-dataset}
\end{figure}

\begin{figure}[htbp]
    \centering
    \includegraphics[width=0.99\textwidth]{diagrams/twooutputs/figure_05_spot_counter_two_output.pdf}
    \caption[Two-output spot counter]{Spot-recognition counts for the two-output models. The plot complements Table~\ref{tab:two_output_models} by showing how often a model recognizes zero, one, or both target spots, rather than only reporting the averaged recognition score.}
    \label{fig:twooutput-spot-counter}
\end{figure}



\begin{table}[htbp]
\centering
\caption{Performance of the two-output models. The best value for each metric is highlighted in bold.}
\label{tab:two_output_models}
\begin{tabular}{p{4.2cm}cccc}
\toprule
Model &
Dice $\uparrow$ &
RMSE / RMSE$_{\mathrm{whole}}$ $\downarrow$ &
Soft Dice$_{\mathrm{whole}}$ $\uparrow$ &
Recognition $\uparrow$ \\
\midrule
Two-output L1 loss, batch size 100
& 0.9222
& 0.0482 / \textbf{0.0096}
& 0.9419
& \textbf{0.9370} \\

Two-output L1 loss
& 0.9156
& 0.0510 / 0.0109
& 0.9327
& 0.9323 \\

Two-output L1 loss, 32 base channels
& 0.8784
& 0.0648 / 0.0146
& 0.9009
& 0.9068 \\

Two-output reconstruction loss
& 0.8800
& \textbf{0.0407} / 0.0099
& \textbf{0.9515}
& 0.9197 \\

Two-output Tversky loss
& \textbf{0.9627}
& 0.6664 / 0.1544
& 0.5341
& 0.6481 \\

Two-output Tversky loss, small dataset
& 0.7293
& 0.5783 / 0.1553
& 0.4926
& 0.5755 \\
\bottomrule
\end{tabular}
\end{table}

Figure~\ref{fig:twooutput-objective-comparison} gives the training-curve context for these models. It shows how the two-output objectives differ during optimization: the L1 and reconstruction objectives preserve intensity more reliably, whereas the Tversky objective emphasizes foreground overlap and therefore reaches a different metric balance.

\subsubsection{Permutation-Invariant L1 Loss}

\begin{figure}[htbp]
    \centering
    \includegraphics[width=0.99\textwidth]{diagrams/twooutputs/twooutput__two-output-L1__bad_02__sample_073451__thesis-comparison.pdf}
    \caption[Two-output L1 failure case]{Failure case for the two-output L1 model. The prediction lacks the structure of the target spot and does not recover a plausible foreground mask.}
    \label{fig:twooutput-l1-bad-example}
\end{figure}


The first two-output baseline uses a permutation-invariant L1 loss on the two predicted spot images. In Table~\ref{tab:two_output_models}, the L1 variants now form the strongest two-output group for recognition and intensity reconstruction. The standard L1 model reaches Dice $0.9156$, \gls{rmse} $0.0510$ / $0.0109$, soft Dice$_{\mathrm{whole}}$ $0.9327$, and recognition $0.9323$. Increasing the batch size to $100$ gives the best recognition among the two-output models, $0.9370$, and the lowest whole-image \gls{rmse}, $0.0096$, while Dice remains high at $0.9222$. The L1 objective therefore no longer appears as a weak separation baseline in the aggregate metrics; instead, it provides a strong intensity-preserving two-output solution, although the selected failure case in Figure~\ref{fig:twooutput-l1-bad-example} shows that individual difficult samples can still fail qualitatively, since the two-output approach is relatively obvious and the L\_1 Loss is very well knows as well as this beeing one of the first training tries. This model will be considered the baseline for future experiments

The additional L1 hyperparameter curves are shown in Figure~\ref{appfig:twooutputmissingloss}. They indicate that the L1 settings converge to competitive validation metrics, with the batch-size-$100$ run slightly ahead of the standard L1 run in Dice, whole-image \gls{rmse}, soft Dice$_{\mathrm{whole}}$, and recognition. The model with only $32$ base channels remains the weakest L1 variant: Dice decreases to $0.8784$, \gls{rmse} rises to $0.0648$ / $0.0146$, soft Dice$_{\mathrm{whole}}$ falls to $0.9009$, and recognition decreases to $0.9068$. This suggests that reduced capacity mainly harms both shape agreement and intensity fidelity, even though the model remains substantially stronger than the older metric values suggested. Representative dense comparison sheets in the appendix material illustrate the remaining qualitative failure modes.

\subsubsection{Composite Reconstruction Loss}

The composite reconstruction loss still performs best on the foreground \gls{rmse} among the two-output models in Table~\ref{tab:two_output_models}, with $0.0407$, and it also gives the best soft Dice$_{\mathrm{whole}}$, $0.9515$. These values show that the reconstruction term is effective at preserving the measured intensity distribution.

Its remaining weakness is relative rather than absolute: compared with the two L1 models, Dice decreases to $0.8800$ and recognition to $0.9197$. Thus, the reconstruction objective is excellent for intensity agreement, but it is slightly less reliable for assigning the two separated objects than the L1 objective. The selected comparison examples support this interpretation and show that the model remains sensitive to dataset morphology and overlap geometry.


\subsubsection{Weighted Tversky Foreground Loss}
\begin{figure}[htbp]
    \centering
    \includegraphics[width=0.95\textwidth]{diagrams/twooutputs/tverky_problem.pdf}
    \caption[Two-output Tversky problem case]{Problem case for the two-output Tversky model. The example shows that the foreground mask can be plausible while the separated intensity assignment remains unstable.}
    \label{fig:twooutput-tversky-problem}
\end{figure}


\begin{figure}[htbp]
    \centering
    \includegraphics[width=0.9\textwidth]{diagrams/twooutputs/tversky/twooutput_tversky_loss_components.pdf}
    \caption[Two-output Tversky loss components]{Loss components for the two-output Tversky model. The mask component converges more reliably than the intensity component.}
    \label{fig:twooutputloss-components}
\end{figure}


The weighted Tversky foreground loss changes the trade-off of the two-output model substantially. It gives the highest Dice in Table~\ref{tab:two_output_models}, $0.9627$, but it is not the strongest model for intensity reconstruction or recognition. Its recognition is only $0.6481$, and its \gls{rmse} values, $0.6664$ / $0.1544$, are much worse than the L1 and reconstruction objectives. The result indicates that foreground-aware supervision produces very strong mask overlap, but in this two-output formulation it does not by itself preserve the separated intensities or consistently recover both spots.

An initial Tversky run was trained on a smaller dataset containing approximately $4{,}000$ samples. This dataset still contained repeated source spots, with individual spots appearing up to three times. The loss curves in Figure~\ref{appfig:loss-small-dataset} show that the training loss decreased, but the model did not reach the performance of the full-dataset Tversky model. As shown in Table~\ref{tab:two_output_models}, the small-dataset model reaches Dice $0.7293$, \gls{rmse} $0.5783$ / $0.1553$, soft Dice$_{\mathrm{whole}}$ $0.4926$, and recognition $0.5755$. These values are clearly below the full-dataset Tversky model for Dice and recognition and remain far below the L1 models for intensity-oriented metrics. The validation data followed the same split strategy as the other experiments and did not contain source spots used during training. The performance gap therefore indicates that the limited dataset size and repeated spot appearances restricted generalization, even though the loss curves themselves appeared stable.

However, the Tversky model has the worst \gls{rmse}, $0.6664$ / $0.1544$, and the lowest soft Dice$_{\mathrm{whole}}$, $0.5341$. In other words, it achieves very strong foreground overlap but its separated intensity values are much less accurate, and this also lowers the recognition score compared with the L1 and reconstruction models. Figure~\ref{fig:twooutputloss-components} shows the same pattern during training: the mask-related component converges more cleanly than the intensity component. The problem case in Figure~\ref{fig:twooutput-tversky-problem} demonstrates the resulting failure mode, where the foreground mask remains plausible while the assigned intensity distribution is unstable. This can be explained by the intensity term having only a weight of $0.2$ in the loss function\footnote{The Tversky loss was later adapted for the multi-output model, where it produced the strongest model in this thesis.}.

The per-dataset results in Figure~\ref{fig:twooutput-per-dataset} are important for interpreting the table. The aggregate values are not caused by uniform performance on all materials. Instead, the models show a strong dependence on the source dataset: easier datasets with compact, well-localized spots dominate the good examples, whereas more difficult datasets with larger deformation, stronger overlap, or less regular spot morphology dominate the failure cases. This pattern is also visible in the selected examples: high-performing cases for the L1 and reconstruction objectives are mostly taken from the IN718 twins dataset, while many of the poor cases come from Ti7Al, Iron, Iron deformed, or Al big grains. The spot-count visualization in Figure~\ref{fig:twooutput-spot-counter} reinforces this point. A model can achieve moderate average spot recognition while still frequently recovering only one of the two spots, or none, on difficult samples. For this reason, the batch-size-$100$ L1 model is the strongest two-output candidate when the priority is reliable spot recognition and intensity-resolved separation, while the Tversky model is mainly useful when foreground overlap is prioritized over intensity fidelity.


\subsection{One-output Model}

The one-output model predicts only one of the two separated spots directly. The second spot is reconstructed as the residual between the input image and the predicted spot. Figure~\ref{fig:oneoutput-objective-comparison} compares the training behavior of the evaluated one-output objectives.

\begin{figure}[htbp]
    \centering
    \includegraphics[width=0.95\textwidth]{diagrams/oneoutput/group_03_one_output_objective.pdf}
    \caption[One-output objective comparison]{Training comparison of the one-output objectives. The Dice/L1 objective gives the stronger aggregate balance between foreground overlap and intensity agreement, whereas the Tversky objective emphasizes mask overlap but does not solve the residual-output limitation in overlap regions.}
    \label{fig:oneoutput-objective-comparison}
\end{figure}

\begin{figure}[htbp]
    \centering
    \includegraphics[width=0.95\textwidth]{diagrams/oneoutput/figure_05_spot_counter_one_output.pdf}
    \caption[One-output spot counter]{Spot-recognition counts for the one-output models. The figure shows how often the residual formulation recovers zero, one, or both separated spots and therefore complements the averaged spot-recognition values in Table~\ref{tab:one_output_models}.}
    \label{fig:oneoutput-spot-counter}
\end{figure}

% One-output models
\begin{table}[htbp]
\centering
\caption{Performance of the one-output models. The best value for each metric is highlighted in bold.}
\label{tab:one_output_models}
\begin{tabular}{p{4.5cm}cccc}
\toprule
Model &
Dice $\uparrow$ &
RMSE / RMSE$_{\mathrm{whole}}$ $\downarrow$ &
Soft Dice$_{\mathrm{whole}}$ $\uparrow$ &
Recognition $\uparrow$ \\
\midrule
One-output Dice/L1 loss
& \textbf{0.6843}
& \textbf{0.1350} / \textbf{0.0337}
& \textbf{0.7997}
& \textbf{0.7747} \\

One-output Tversky loss
& 0.6291
& 0.1652 / 0.0390
& 0.6314
& 0.7319 \\
\bottomrule
\end{tabular}
\end{table}

Table~\ref{tab:one_output_models} shows that the Dice/L1 model is the stronger one-output model overall. It reaches Dice $0.6843$, \gls{rmse} $0.1350$ / $0.0337$, soft Dice$_{\mathrm{whole}}$ $0.7997$, and recognition $0.7747$. The per-dataset performance differs substantially as well, as shown in Figure~\ref{fig:oneoutput-dataset-heatmap}. This wide spread is the central weakness of the one-output formulation.

\begin{figure}[htbp]
    \centering
    \includegraphics[width=0.95\textwidth]{diagrams/oneoutput/oneoutput_problem.pdf}
    \caption[One-output problem case]{Problem case for the one-output model. The residual construction cannot preserve shared overlap intensity reliably, which leads to missing or distorted signal in the second spot.}
    \label{fig:oneoutputpreviewexample}
\end{figure}

The main limitation of this model is shown in Figure~\ref{fig:oneoutputpreviewexample}. The directly predicted spot captures parts of the target morphology, but its spatial extent is underestimated. The second output, constructed as the residual image, deviates more strongly from the target because intensity removed for the first spot is no longer available for the second one. This is especially problematic in the overlap region, where the measured intensity is the sum of both physical spots. The network would need to assign shared intensity to both separated outputs, but the residual formulation only allows the intensity to remain in one channel.

The Tversky model changes the metric profile, but it still does not truly separate the two spots. Instead, the qualitative examples in the supplementary handout indicate that it tends to crop out one plausible spot while leaving the residual output incomplete. Compared with the Dice/L1 model, Dice decreases from $0.6843$ to $0.6291$, \gls{rmse} worsens from $0.1350$ / $0.0337$ to $0.1652$ / $0.0390$, soft Dice$_{\mathrm{whole}}$ drops from $0.7997$ to $0.6314$, and recognition decreases from $0.7747$ to $0.7319$. Thus, the Tversky loss sacrifices intensity fidelity and does not make the residual representation more robust in the aggregate metrics.

The spot-counter plot in Figure~\ref{fig:oneoutput-spot-counter} is important for interpreting these averages. A moderate spot-recognition value does not mean that both spots are recovered consistently. Instead, the model often predicts one plausible spot while the residual channel remains incomplete or distorted. This explains why the one-output Dice/L1 model can look competitive in individual good examples while still failing on difficult samples.


\begin{figure}[htbp]
    \centering
    \includegraphics[width=1\textwidth]{diagrams/oneoutput/figure_04_dataset_model_heatmap_one_output.pdf}
    \caption[One-output per-dataset metrics]{Per-dataset metrics for the one-output models. The heatmap shows that the residual formulation is particularly unstable on difficult datasets with strong overlap, deformation, or irregular spot morphology.}
    \label{fig:oneoutput-dataset-heatmap}
\end{figure}

The per-dataset heatmap in Figure~\ref{fig:oneoutput-dataset-heatmap} shows that this behavior is not uniform across the data. The best selected Dice/L1 cases are concentrated in the IN718 twins dataset, where the model can reconstruct compact examples with high Dice and spot recognition. The weak Dice/L1 cases are mainly from Ti7Al and Iron. For the Tversky model, strong examples occur across Ti7Al, Al, IN718 twins, Al big grains, and Al small grains, but the worst cases include IN718 twins, Al deformed LoG, Ti7Al, and Iron. This distribution indicates that the residual formulation is sensitive to both overlap geometry and source-dataset morphology. It can work on favorable cases, but it does not generalize reliably when the two spot contributions are strongly entangled.

Overall, the one-output approach is useful as a baseline but remains physically limited. The Dice/L1 objective gives the better balance of spot recognition and intensity accuracy, while the Tversky objective worsens all metrics shown in Table~\ref{tab:one_output_models}. Compared with the two-output models, the one-output formulation avoids channel-assignment ambiguity, but it introduces a stronger representational bottleneck: the second spot is not independently predicted and therefore cannot retain intensity that has already been assigned to the first output.


\subsection{Multi-output Model}

% Multi-output models
\begin{table}[htbp]
\centering
\caption{Performance of the multi-output models. The best value for each metric is highlighted in bold.}
\label{tab:multi_output_models}
\resizebox{\textwidth}{!}{%
\begin{tabular}{p{5.2cm}cccc}
\toprule
Model &
Dice $\uparrow$ &
RMSE / RMSE$_{\mathrm{whole}}$ $\downarrow$ &
Soft Dice$_{\mathrm{whole}}$ $\uparrow$ &
Recognition $\uparrow$ \\
\midrule
Multi-output Tversky loss, ReduceLROnPlateau factor 0.5
& \textbf{0.9411}
& \textbf{0.0396} / \textbf{0.0084}
& \textbf{0.9758}
& \textbf{0.9507} \\

Multi-output Tversky loss, learning rate $10^{-4}$
& \textbf{0.9411}
& 0.0432 / 0.0104
& 0.9728
& 0.9490 \\

Multi-output Tversky loss, batch size 20
& 0.9380
& 0.0401 / 0.0096
& 0.9729
& 0.9490 \\

Multi-output Tversky loss, learning rate $5\times10^{-5}$
& 0.9349
& 0.0409 / 0.0094
& 0.9686
& 0.9470 \\

Multi-output Tversky loss, batch size 100
& 0.9332
& 0.0411 / 0.0095
& 0.9704
& 0.9460 \\

Multi-output Tversky loss, 32 base channels
& 0.9294
& 0.0431 / 0.0104
& 0.9685
& 0.9432 \\

Multi-output Tversky loss, learning rate $3\times10^{-4}$
& 0.9086
& 0.0492 / 0.0119
& 0.9700
& 0.9297 \\

Multi-output BCE/MSE loss
& 0.8214
& 0.1273 / 0.0309
& 0.7869
& 0.8471 \\

Multi-output BCE/MSE loss, small dataset
& 0.6439
& 0.1460 / 0.0497
& 0.6283
& 0.7490 \\

\bottomrule
\end{tabular}%
}
\end{table}

\begin{figure}[htbp]
    \centering
    \includegraphics[width=0.95\textwidth]{diagrams/multihead/group_04_multi_head_objective_scheduler.pdf}
    \caption[Multi-output objective and scheduler comparison]{Comparison of the multi-output objectives and scheduler settings. The binary cross-entropy/mean squared error baseline is clearly weaker than the Tversky models, while the Tversky model with ReduceLROnPlateau gives the strongest balance of segmentation quality, intensity reconstruction, and spot recognition.}
    \label{fig:multihead-objective-scheduler}
\end{figure}

\subsubsection{Binary Cross-Entropy and Mean Squared Error Loss}

Table~\ref{tab:multi_output_models} shows that the full-dataset \gls{bce}/\gls{mse} model is stronger than the one-output baselines, but remains substantially weaker than the Tversky multi-output models. It reaches Dice $0.8214$, \gls{rmse} $0.1273$ / $0.0309$, soft Dice$_{\mathrm{whole}}$ $0.7869$, and recognition $0.8471$. The selected examples show that this average hides a considerable lower tail: strong \gls{bce}/\gls{mse} cases can be very accurate, while poor cases still fail to recover both separated spots reliably.

The \gls{bce}/\gls{mse} loss was also tested on the smaller dataset. The loss curves in Figure~\ref{appfig:loss-small-dataset} indicate overfitting: the training loss continued to converge, while the validation loss worsened. The model was nevertheless evaluated on unseen validation data. Its metrics are lower than those of the full-dataset \gls{bce}/\gls{mse} model, with Dice decreasing from $0.8214$ to $0.6439$, \gls{rmse} worsening from $0.1273$ / $0.0309$ to $0.1460$ / $0.0497$, soft Dice$_{\mathrm{whole}}$ decreasing from $0.7869$ to $0.6283$, and recognition decreasing from $0.8471$ to $0.7490$. This confirms that the small dataset limits generalization across all metrics shown in the main table. Nevertheless, even this small-dataset multi-output baseline remains competitive with, and in several metrics stronger than, the one-output approaches.

\subsubsection{Tversky Loss and Hyperparameter Tuning}

Switching to the Tversky foreground loss produces the decisive improvement. The selected scheduler configuration in Table~\ref{tab:multi_output_models} reaches Dice $0.9411$, \gls{rmse} $0.0396$ / $0.0084$, soft Dice$_{\mathrm{whole}}$ $0.9758$, and recognition $0.9507$ on the validation data. Relative to the full-dataset \gls{bce}/\gls{mse} baseline, this improves Dice by $0.1197$, foreground \gls{rmse} by $0.0877$, whole-image \gls{rmse} by $0.0225$, soft Dice$_{\mathrm{whole}}$ by $0.1889$, and recognition by $0.1036$.

The hyperparameter-tuning results in Table~\ref{tab:multi_output_models} show that the Tversky loss is robust, but still sensitive to capacity and optimization settings. Reducing the model to $32$ base channels lowers Dice to $0.9294$ and recognition to $0.9432$, and increases \gls{rmse} to $0.0431$ / $0.0104$. The problem example in Figure~\ref{fig:multihead-base32-problem} illustrates the practical effect of this reduction: the smaller model can still solve many samples, but it is more likely to lose structure on difficult overlaps. Increasing the batch size to $100$ also underperforms the best settings, with Dice $0.9332$, \gls{rmse} $0.0411$ / $0.0095$, and recognition $0.9460$. The batch-size-$20$ run is stronger, reaching Dice $0.9380$ and recognition $0.9490$, which suggests that the optimization benefits from a moderately smaller batch.

\begin{figure}[htbp]
    \centering
    \includegraphics[width=0.95\textwidth]{diagrams/multihead/tversky_base_32_problem.pdf}
    \caption[Multi-output base channel problem example]{Problem case for the smaller base channel multi-output Tversky model. The example illustrates that reducing model capacity increases the risk of failures on difficult overlaps.}
    \label{fig:multihead-base32-problem}
\end{figure}

The learning rate tuning shows a similar pattern. A learning rate of $3\times10^{-4}$ is the weakest Tversky setting in the main table, with Dice $0.9086$, \gls{rmse} $0.0492$ / $0.0119$, and recognition $0.9297$. The corresponding loss curves in Figure~\ref{fig:multihead-learning-rate-comparison} show that this run does not converge satisfactorily: the optimization remains unstable compared with the other learning rate settings and does not reach the same validation quality. The lower learning rate of $5\times10^{-5}$ is considerably stronger, but remains below the selected scheduler run. The fixed learning rate of $10^{-4}$ ties the scheduler run on Dice at $0.9411$, but the scheduler gives better \gls{rmse}, soft Dice$_{\mathrm{whole}}$, and recognition. Therefore, the scheduler model was selected as the final model. The learning rate behavior is shown explicitly in Figure~\ref{fig:multihead-tversky-learning-rate}, and the loss components in Figure~\ref{fig:multihead-tversky-validation-components} show that the selected model keeps both mask and intensity terms stable during training.

\begin{figure}[htbp]
    \centering
    \includegraphics[width=0.95\textwidth]{diagrams/multihead/multi_head_tversky_plateau_validation_components.pdf}
    \caption[Multi-output Tversky validation components]{Validation components for the multi-output Tversky model with ReduceLROnPlateau. The component curves show how the final model balances mask and intensity quality.}
    \label{fig:multihead-tversky-validation-components}
\end{figure}

\begin{figure}[htbp]
    \centering
    \includegraphics[width=0.95\textwidth]{diagrams/multihead/multi_head_tversky_plateau_learning_rate.pdf}
    \caption[Multi-output Tversky learning rate schedule]{Learning rate schedule for the multi-output Tversky model with ReduceLROnPlateau. The learning rate is reduced by a factor of $0.5$ after plateaus of $30$ epochs.}
    \label{fig:multihead-tversky-learning-rate}
\end{figure}

\begin{figure}[htbp]
    \centering
    \includegraphics[width=0.95\textwidth]{diagrams/multihead/group_05_multi_head_learning_rate.pdf}
    \caption[Multi-output learning rate comparison]{Comparison of learning rate settings for the multi-output Tversky models. The fixed learning rate runs and the scheduler run show that overly small learning rates underfit, while the scheduler gives the best balance of spot recognition and intensity error.}
    \label{fig:multihead-learning-rate-comparison}
\end{figure}

\begin{figure}[htbp]
    \centering
    \includegraphics[width=0.95\textwidth]{diagrams/multihead/group_06_multi_head_batch_base.pdf}
    \caption[Multi-output batch and base channel comparison]{Comparison of batch-size and base channel settings for the multi-output Tversky models. Reducing the base channels and using a batch size of $100$ both reduce segmentation quality compared with the stronger Tversky settings.}
    \label{fig:multihead-batch-base-comparison}
\end{figure}

Figures~\ref{fig:multihead-learning-rate-comparison} and~\ref{fig:multihead-batch-base-comparison} summarize the hyperparameter-tuning process visually. Together with Table~\ref{tab:multi_output_models}, they show that no single change explains the final performance. The per-dataset heatmap in Figure~\ref{fig:multihead-dataset-heatmap} further shows that the multi-output Tversky models improve the lower-performing material groups compared with the \gls{bce}/\gls{mse} baseline, but that the task remains dataset dependent. Difficult material groups and strongly overlapping examples still form the lower tail, so the aggregate score should be interpreted together with the dataset-wise behavior. The final model is therefore the result of combining the multi-output architecture, the Tversky foreground objective, sufficient model capacity, a moderate batch size, and an adaptive learning rate schedule.

\begin{figure}[htbp]
    \centering
    \includegraphics[width=0.95\textwidth]{diagrams/multihead/example_best_model.pdf}
    \caption[Best multi-output model example]{Representative example for the selected multi-output model. The prediction keeps both mask localization and intensity reconstruction substantially more stable than the one-output and two-output alternatives.}
    \label{fig:multihead-best-example}
\end{figure}

\subsection{Final Model Evaluation on the Test Dataset}

Before applying the final model to the held-out test dataset, the most relevant validation candidates were compared across the output formulations. Table~\ref{tab:final_model_candidate_comparison} summarizes this selection step. The one-output Dice/L1 model remains physically limited by its residual construction, and the two-output L1 model with batch size $100$ is strong but still couples localization and intensity reconstruction in the same output image. The multi-output \gls{bce}/\gls{mse} baseline improves on the one-output formulation, but it is clearly weaker than the multi-output Tversky scheduler model. The ReduceLROnPlateau configuration is therefore selected as the final model because it reaches the best \gls{rmse}, best soft Dice$_{\mathrm{whole}}$, and best recognition while also achieving the highest Dice together with the fixed $10^{-4}$ run.

\begin{table}[htbp]
\centering
\caption{Validation comparison of the most relevant model candidates.}
\label{tab:final_model_candidate_comparison}
\begin{tabular}{p{5.2cm}cccc}
\toprule
Model &
Dice $\uparrow$ &
RMSE / RMSE$_{\mathrm{whole}}$ $\downarrow$ &
Soft Dice$_{\mathrm{whole}}$ $\uparrow$ &
Recognition $\uparrow$ \\
\midrule
One-output Dice/L1 loss
& 0.6843
& 0.1350 / 0.0337
& 0.7997
& 0.7747 \\

Two-output L1 loss, batch size 100
& 0.9222
& 0.0482 / 0.0096
& 0.9419
& 0.9370 \\

Multi-output BCE/MSE loss
& 0.8214
& 0.1273 / 0.0309
& 0.7869
& 0.8471 \\

Multi-output Tversky loss, ReduceLROnPlateau factor 0.5
& \textbf{0.9411}
& \textbf{0.0396} / \textbf{0.0084}
& \textbf{0.9758}
& \textbf{0.9507} \\
\bottomrule
\end{tabular}
\end{table}

The corresponding per-dataset heatmaps in Figures~\ref{fig:oneoutput-dataset-heatmap},~\ref{fig:twooutput-per-dataset}, and~\ref{fig:multihead-dataset-heatmap} show that the candidate ranking is not uniform across all source datasets. Across the compared models, the following dataset-level tendencies are visible:
\begin{itemize}
    \item \texttt{IN718\_twins} is the most consistently favourable dataset. It contributes many of the strongest one-output and two-output examples and remains well handled by the multi-output models, suggesting that compact, well-localized spot morphologies are easier to separate.
    \item The aluminium datasets with smaller or regular grains are generally more stable than the difficult material groups. Their performance improves strongly with the multi-output Tversky model, indicating that the final architecture can use the additional mask and intensity heads to reduce dataset-specific failures.
    \item \texttt{Ti7Al}, \texttt{Iron}, \texttt{Iron\_deformed}, and partly \texttt{Al\_big\_grains} form the lower-performing group in the one-output and two-output comparisons. These datasets contain more difficult morphology or overlap geometry, so the models more often recover only one spot or produce less accurate intensity separation.
    \item An interesting pattern is that \texttt{IN718\_twins} appears in both strong and weak examples. Dataset identity alone therefore does not explain the result; the overlap geometry and individual spot morphology still matter within a nominally favourable dataset.
    \item The multi-output Tversky model reduces the gap between easy and difficult datasets most clearly. The remaining weak cases are no longer a failure of one specific material only, but are concentrated in high-overlap or irregular samples across the harder groups.
\end{itemize}

\begin{table}[htbp]
\centering
\caption{Performance of the selected final model on the held-out test dataset.}
\label{tab:final_model_test_evaluation}
\begin{tabular}{p{5.2cm}cccc}
\toprule
Model &
Dice $\uparrow$ &
RMSE / RMSE$_{\mathrm{whole}}$ $\downarrow$ &
Soft Dice$_{\mathrm{whole}}$ $\uparrow$ &
Recognition $\uparrow$ \\
\midrule
Multi-output Tversky loss, ReduceLROnPlateau factor 0.5 evaluation
& 0.9300
& 0.0460 / 0.0105
& 0.9581
& 0.9398 \\
\bottomrule
\end{tabular}
\end{table}

The selected model was then evaluated on a held-out test dataset of $20{,}000$ synthetically generated overlaps. Table~\ref{tab:final_model_test_evaluation} reports Dice $0.9300$, \gls{rmse} $0.0460$ / $0.0105$, soft Dice$_{\mathrm{whole}}$ $0.9581$, and recognition $0.9398$. These values are lower than the validation metrics of the selected scheduler run, but they remain close enough to indicate that the model generalizes well to unseen generated samples. The representative prediction in Figure~\ref{fig:multihead-best-example} shows the main advantage of the final model: the two masks remain spatially separated while the corresponding intensity heads preserve the spot signal.

\begin{figure}[htbp]
    \centering
    \includegraphics[width=0.95\textwidth]{diagrams/multihead/figure_04_dataset_model_heatmap_multi_output.pdf}
    \caption[Multi-output per-dataset metrics]{Per-dataset metric heatmap for the multi-output models. The heatmap shows that the multi-output Tversky models improve performance across datasets, while difficult material datasets still form the lower tail.}
    \label{fig:multihead-dataset-heatmap}
\end{figure}

\begin{figure}[htbp]
    \centering
    \includegraphics[width=0.95\textwidth]{diagrams/multihead/figure_05_spot_counter_multi_output.pdf}
    \caption[Multi-output spot-count performance]{Spot-count grouped performance for the multi-output models. The final Tversky scheduler model recovers both spots more consistently than the binary cross-entropy/mean squared error baseline and the reduced-capacity models.}
    \label{fig:multihead-spot-counter}
\end{figure}

The dataset heatmap in Figure~\ref{fig:multihead-dataset-heatmap} and the spot-count plot in Figure~\ref{fig:multihead-spot-counter} show that the final model is not perfect, but its errors are much less severe than in the previous output models. The best selected cases reach very high recognition and Dice values with very low \gls{rmse}, while the remaining weak cases are concentrated in high-overlap or difficult material examples. Nevertheless, the multi-output Tversky model with ReduceLROnPlateau is the strongest evaluated model because it combines accurate foreground separation with substantially improved intensity reconstruction and a much smaller failure tail.

